% TOP_PROBLEM: {"id": 30002424, "title": "Gottschalk Surjunctivity Problem", "category_id": 11, "category": {"id": 11, "name": "dynamical_systems", "display_name": "Dynamical Systems", "description": "Problems about long-term behavior of deterministic systems, Hamiltonian dynamics, and periodic orbits.", "slug": "dynamical-systems", "order_index": 11, "created_at": "2026-07-31T15:26:25.671Z"}, "status": "open"}
% FIRST_POSTED: 2026-09-25
% ENTRY_KIND: statement_only
% !TeX program = lualatex
% Definition and sources only; this file is not an LLM solution attempt.
\documentclass[11pt]{article}
\usepackage[margin=25mm]{geometry}
\usepackage{fontspec}
\setmainfont{Latin Modern Roman}
\usepackage{amsmath,amssymb,mathtools}
\usepackage{unicode-math}
\setmathfont{Latin Modern Math}
\usepackage{xurl,hyperref}
\hypersetup{colorlinks=true,urlcolor=blue}
\setcounter{secnumdepth}{0}
\setlength{\parindent}{0pt}
\setlength{\parskip}{5pt}
\setlength{\emergencystretch}{3em}
\begin{document}
\section*{\texorpdfstring{Gottschalk Surjunctivity Problem}{Gottschalk Surjunctivity Problem}}\label{30002424}
\subsection{Definitions and mathematical statement}
Let $G$ be any group and $A$ a finite alphabet. Give $A^G$ the product topology with $A$ discrete, and let $G$ act by shifts $(g\cdot x)(h)=x(g^{-1}h)$. A cellular automaton is a continuous shift-equivariant map $F:A^G\to A^G$, equivalently a finite-memory rule: for some finite $D\subset G$ and $\mu:A^D\to A$,
\[
 F(x)(g)=\mu\bigl((x(gh))_{h\in D}\bigr).
\]
Surjunctivity asserts
\[
 F\text{ injective}\quad\Longrightarrow\quad F\text{ surjective}
\]
for every $G,A,F$. The local rule is uniform over all sites. Results about non-uniform automata must therefore be related back to this uniform target rather than changing its definition. Neither finite generation nor countability of $G$ is imposed in the record.
\par\smallskip\textit{Formulation and concept references:} \href{https://arxiv.org/html/2503.23435v2}{[S1]}.

\subsection{Short English statement}
On every group, must a finite-alphabet cellular automaton that never merges two configurations also reach every configuration?
\subsection{Sources}
\begin{itemize}
\item[S1]Xuan Kien Phung, On Gottschalk's surjunctivity conjecture for non-uniform cellular automata, arXiv:2503.23435v2, 2026.
\url{https://arxiv.org/html/2503.23435v2}.
\textit{Formulation source.}
\end{itemize}
\end{document}
